% !TeX root = RJwrapper.tex
\title{Kernopt: A Package for Estimating Count Data Distributions with a Discrete Symmetric Optimal Kernel.}


\author{by Gilles Durrieu, Thomas Fillon, and Tristan Senga Kiessé}

\maketitle

\abstract{%
\CRANpkg{kernopt} is an R package that implements a discrete symmetric optimal kernel for estimating count data distributions. The nonparametric estimator using the discrete symmetric optimal kernel was illustrated on simulated and real-word data included in the package, in comparison with two other discrete symmetric kernels.
}

\makeatletter
\newsavebox\pandoc@box
\newcommand*\pandocbounded[1]{% scales image to fit in text area
  \sbox\pandoc@box{#1}%
  \ifdim\wd\pandoc@box>\linewidth
    \resizebox{\linewidth}{!}{\copy\pandoc@box}%
  \else
    \copy\pandoc@box
  \fi
}
\makeatother

\section{Introduction}\label{introduction}

Since the original works of Rosenblatt \citep{Rosenblatt1956} and Parzen \citep{Parzen1962}, nonparametric kernel smoothing was often studied for estimating continuous data distributions by using symmetric and asymmetric continuous kernels, e.g., Gaussian, Epanechnikov \citep{Epan69}, on the one hand, beta \citep{Chen99}, inverse Gaussian \citep{scaillet2004}, on the other hand, respectively. Then, discrete kernel smoothing was developed to specifically address the nonparametric estimation of count data distributions by using discrete kernels associated to discrete probability distributions, such as the kernel of \citet{Racine2020} for ordered data, Conway-Maxwell-Poisson kernel \citep{Huang2022} and, more recently, underdispersed count kernels \citep{Kokon_etal_2023}. That was a valuable alternative to the traditional parametric Poisson modeling, which is often of limited use because empirical count data sets typically exhibit over/under-dispersion and/or an excess number of zeros \citep{KokonSenga11}. One R package implements the kernel method for modeling distributions of continuous (bounded and positive) and discrete (categorical and count) data (\CRANpkg{Ake} package by \citet{Wansouwe2016}). Nevertheless, kernel density estimation using traditional continuous symmetric kernels (e.g., Gaussian, Epanechnikov, rectangular, triangular) is available by using the function \texttt{density} available in \texttt{R}. Various functions for kernel smoothing from the works of \citet{WandJones1995} were also implemented in \CRANpkg{KernSmooth} package \citep{Kernsmooth}. In addition, multivariate kernel density estimation was developed in \CRANpkg{ks} package \citep{Duong2007}.

The wide variety of discrete associated kernels available in the literature raised the issue of finding one discrete ``optimal'' associated kernel, which was developed by \citet{SengaDurrieu2024}. That new discrete symmetric ``optimal'' kernel is implemented in the R package \CRANpkg{kernopt} which is also available at \url{https://github.com/thomasfillon/kernopt}. Other discrete associated kernels are also available, e.g., Epanechnikov \citep{chu2017}, triangular and binomial \citep{Wansouwe2016}, for a comparison with the new discrete kernel developed. In addition to the choice of discrete kernel, the choice of the bandwidth parameter is challenging. We propose a cross-validation method for selecting the smoothing bandwidth to be used with the discrete kernel. Illustrations are provided on discrete simulated and real-word application data.

One can refer to the article of \citet{SengaDurrieu2024} for technical details on mathematical results (propositions, theorem and proofs) provided in that paper.
The \CRANpkg{kernopt} package not only implements the discrete symmetric optimal kernel but also introduces a data-driven bandwidth selection method based on cross-validation, enhancing the estimator's adaptability and performance to various data.

By providing a comprehensive and user-friendly implementation, the \CRANpkg{kernopt} package plays a pivotal role in enabling researchers and practitioners to effectively apply nonparametric optimal kernel estimation to discrete count data. Available on both CRAN and GitHub, the package ensures broad accessibility and facilitates reproducibility in statistical analyses. Example usages of the package on both simulated data and real aquaculture data from the SIMTAP project demonstrate the better estimation performance of our proposed method.

The paper is organized as follows. \hyperref[sec2-nonparametric-kernel-estimation]{Section 2} presents the fundamental principles of the nonparametric estimation method using discrete associated kernels for probability mass functions. In \hyperref[sec3-optimal-kernel]{Section 3}, we describe the new discrete symmetric ``optimal'' kernel developed by \citet{SengaDurrieu2024}, with the procedure for bandwidth selection. The performance of the estimator with the discrete symmetric ``optimal'' kernel is illustrated in \hyperref[sec4-presentation-simu]{Section 4} on simulated data, then applications on real data are provided in \hyperref[sec5-application-real-data]{Section 5}.

\section{Nonparametric discrete kernel density estimator}\label{sec2-nonparametric-kernel-estimation}

Let \(X_{1}, \ldots, X_{n}\) be a sequence of independent and identically distributed discrete random variable having an unknown probability mass function (pmf) \(f\) to estimate on discrete support \(\mathcal{S}\subseteq \mathbb{Z}\), with \(\mathbb{Z}\) being the set of all integer numbers.
For a target point \(x\in\mathcal{S}\) and a bandwidth parameter \(h>0\), the nonparametric estimator \(\widehat{f}_{K,h}\) of \(f\) is given by

\begin{equation}
\widehat{f}_{K,h}(x)=\frac{1}{n}\sum_{i=1}^{n}K_{x,h}(X_{i}),
\label{eq:estim-fn}
\end{equation}
where \(K_{x,h}\) is a symmetric pmf on discrete support \(\mathcal{S}_{x}\), i.e.~\(\sum_{y\in\mathcal{S}_{x}}K_{x,h}(y)=1\), \(0\leq K_{x,h}(y)\leq 1\) and \(K_{x,h}(y)=K_{x,h}(-y)\). More precisely, a discrete associated kernel \(K_{x,h}(\cdot)\) with a support \(\mathcal{S}_{x}\) (which contains \(x\)) is a pmf of random variable \(\mathcal{K}_{x,h}\), whose modal probability, mean and variance behave asymptotically as follows when \(h\to 0\):
\begin{equation}
K_{x,h}(x)=\Pr(\mathcal{K}_{x,h}=x) \to 1, \text{ \ }\mathbb{E}(\mathcal{K}_{x,h}) \to x \text{ \ and\ } \text{Var}(\mathcal{K}_{x,h}) \to 0.
\label{eq:ModalProb-mean-var-Kxh}
\end{equation}

The discrete kernel estimator \eqref{eq:estim-fn} is defined up to the normalizing constant \(C_{n}=\sum_{x\in\mathcal{S}}\widehat{f}_{K,h}(x)\) in \citet{Senga2017}.

We particularly considered a symmetric pmf \(K_{x,h}\) satisfying the conditions in \eqref{eq:ModalProb-mean-var-Kxh} such that
\begin{equation} 
 \mathbb{E}(\mathcal{K}_{x,h}) =x, \text{ \ }\mathbb{E}(\mathcal{K}^{2}_{x,h}) =\tau^{2}(x,h)<\infty \text{ and } \text{Var}(\mathcal{K}_{x,h})=\tau^{2}(x,h)- x^{2}<\infty,  
\label{eq:condition-symm-kern}
\end{equation}
where \(\text{Var}(\mathcal{K}_{x,h})\) does not depend on \(x\) and tends to \(0\) as \(h\) goes to \(0\).

\subsection{Choice of discrete kernels}\label{choice-of-discrete-kernels}

We consider two examples of discrete symmetric kernels with mean value equal to \(x\), i.e., \(\mathbb{E}(\mathcal{K}_{x,h}) =x\).

For a positive integer \(h\in\mathbb{N}\setminus\{0\}\), the first example is a discrete symmetric Epanechnikov kernel with \(\mathcal{S}_{x}=\{x-h,\ldots,x,\ldots,x+h\}\) given by

\begin{equation*}  
K^{Epan}_{x,h}(y)= \left(a\bigg(\frac{x-y}{h}\bigg)^{2} + b\right){\bf 1}_{\{x-h,\ldots,x,\ldots,x+h\}}(y),
\end{equation*}
with \(a=-b=3h/(1-4h^{2})\); see \citet{chu2017}. That kernel only involves the parameter \(h\), which is a positive integer that scales the distance between \(x\) and \(y\). The discrete Epanechnikov kernel is less often used in the literature, compared to its continuous counterpart.

\subsubsection{\texorpdfstring{\emph{Proposition 1}}{Proposition 1}}\label{proposition-1}

\emph{
The random variable \(\mathcal{K}^{Epan}_{x,h}\) associated to the discrete Epanechnikov kernel with support \(\mathcal{S}_{x}=\{x-h,\ldots,x,\ldots,x+h\}\), with \(h\in\mathbb{N}\setminus\{0\}\), has a Dirac distribution at \(x\) as \(h\to 1\) and a discrete uniform distribution as \(h\to \infty\).
}\\

The \CRANpkg{kernopt} R package computes that kernel at a target \(x\) for various values of observations \(z\) and a fixed bandwidth \(h\) by using the function \texttt{discrete\ kernel} such that (Figure \ref{fig:kernopt-epanech}):

\begin{verbatim}
x <- 5 # Target
z <- 0:10 # Observations
h <- c(1, 2, 3, 4) # Set of bandwidths
K_epan <- matrix(
  data = 0,
  nrow = length(z),
  ncol = length(h)
)
for (i in 1:length(h))
{
  K_epan[, i] <- discrete_kernel(kernel = "epanech", x, z, h[i])
}
plot(
  z,
  K_epan[, 1],
  xlab = "x",
  ylab = "Probability",
  ylim = c(0, 1),
  pch = 1
)
lines(z, K_epan[, 1], lty = 1)
for (i in 2:length(h))
{
  points(z, K_epan[, i], xlab = "z", pch = i)
  lines(z, K_epan[, i], lty = i)
}
legend(
  "topleft",
  c("h=1", "h=2", "h=3", "h=4"),
  lty = 1:4,
  pch = 1:4,
  cex = 1.6
)
\end{verbatim}

\begin{figure}
\centering
\pandocbounded{\includegraphics[keepaspectratio,alt={\label{fig:kernopt-epanech}Distribution of discrete symmetric Epanechnikov kernel at the target point x=5 for various bandwidth parameters h.}]{RJ-2026-018_files/RJ-2026-018_files/figure-latex/kernopt-epanech-1.pdf}}
\caption{\label{fig:kernopt-epanech}Distribution of discrete symmetric Epanechnikov kernel at the target point \(x=5\) for various bandwidth parameters \(h\).}
\end{figure}

For \(h>0\) and non-negative integer \(a\in\mathbb{N}\), the second example is the discrete symmetric triangular kernel with \(\mathcal{S}_{x}=\{x-a,\ldots,x,...,x+a\}\), which has a pmf defined by

\begin{equation*}
K^{Triang}_{a;x,h}(y)= \frac{(a+1)^{h}}{C(a,h)}\left(1 - \frac{|y-x|^{h}}{(a+1)^{h}}\right){\bf 1}_{\{x-a,\ldots,x,\ldots,x+a\}}(y),
\end{equation*}

with \(C(a,h)=(2a+1)(a+1)^{h} - 2\sum_{j=0}^{a}j^{h}\), see \citet{KokonZocSenga2007}. That kernel depends on two parameters \(h\) and \(a\). That kernel was implemented in the \CRANpkg{kernopt} package through the function \texttt{discrete\_kernel} with its specific arguments as follows (Figure \ref{fig:kernopt-triangle}):

\begin{verbatim}
x <- 5
z <- 0:10
h <- c(0.1, 0.4, 1, 2)
a <- 1

K_trg <- matrix(
  data = 0,
  nrow = length(z),
  ncol = length(h)
)

for (i in 1:length(h))
{
  K_trg[, i] <- discrete_kernel(kernel = "triang", x, z, h[i], a)
}
plot(
  z,
  K_trg[, 1],
  xlab = "x",
  ylab = "Probability",
  ylim = c(0, 1),
  pch = 1
)
lines(z, K_trg[, 1], lty = 1)
for (i in 2:length(h))
{
  points(z, K_trg[, i], xlab = "z", pch = i)
  lines(z, K_trg[, i], lty = i)
}
legend(
  "topleft",
  c("h=0.1", "h=0.4", "h=1", "h=2"),
  lty = 1:4,
  pch = 1:4,
  cex = 1.6
)
\end{verbatim}

\begin{figure}
\centering
\pandocbounded{\includegraphics[keepaspectratio,alt={\label{fig:kernopt-triangle}Distribution of discrete symmetric triangular kernel (a=1) at the target point x=5 for various bandwidth parameters h.}]{RJ-2026-018_files/RJ-2026-018_files/figure-latex/kernopt-triangle-1.pdf}}
\caption{\label{fig:kernopt-triangle}Distribution of discrete symmetric triangular kernel (\(a=1\)) at the target point \(x=5\) for various bandwidth parameters \(h\).}
\end{figure}

The discrete symmetric triangular kernel is also available in the \CRANpkg{Ake} R package \citep{Wansouwe2016}.

Other discrete kernels are available in the literature such as the Habbema kernel \citep{Tutz1996}, which is implemented in the \CRANpkg{NPHazardRate} R package \citep{Bagkavos2018} for estimating discrete data of time failure.

\subsubsection*{Mean Summed Squared Error}\label{mean-summed-squared-error}
\addcontentsline{toc}{subsubsection}{Mean Summed Squared Error}

For \(x\in \mathcal{S}\) and \(h>0\), the global performance of the estimator \(\widehat{f}_{K,h}\) of pmf \(f\) can be assessed by calculating its Mean Summed Squared Error (MSSE) given by
\begin{eqnarray}
 \text{MSSE}\left(\widehat{f}_{K,h}(x)\right)
 &=&\frac{1}{n}\sum_{x\in\mathcal{S}}f(x)\sum_{y\in\mathcal{S}_{x}}K^{2}_{x,h}(y) + \sum_{x\in\mathcal{S}}\bigg(\frac{f^{(2)}(x)}{2}\mathrm{Var}(\mathcal{K}_{x,h}) \bigg)^{2} + o\bigg(\frac{1}{n}+h^{2}\bigg)  \label{eq:mise}
 \end{eqnarray}
where the finite difference \(f^{(2)}\) of order \(2\) of \(f\) is such that \(f^{(k)}(x)=\{f^{(k-1)}(x)\}^{(1)}\) with
\begin{equation*}
 f^{(1)}(x)=\left\{   \begin{array}{cc}
      \left(f(x+1)-f(x-1)\right)/2,   & \mathrm{if \ } x\in\mathbb{N}\setminus\{0\},  \\
      f(1)- f(0),   & \mathrm{if \ } x=0 
    \end{array}
    \right.
\end{equation*}
and
\begin{equation*}
 f^{(2)}(x)=\left\{   
 \begin{array}{cc}
\left(f(x+2)-2f(x) + f(x-2)\right)/4,& \mathrm{if \ } x\in\mathbb{N}\setminus\{0, 1\},\\
\left(f(3) - 3f(1) + f(0)\right)/4,& \mathrm{if \ } x=1, \\
\left(f(2) - 2f(1) + f(0)\right)/2& \mathrm{if \ } x=0 .
    \end{array}
    \right.
\end{equation*}
At the end of \hyperref[sec3-optimal-kernel]{Section 3} we provide detailed examples of the expression of MSSE of \(\widehat{f}_{K,h}\) of the pmf \(f\) as a function of the discrete symmetric kernel considered.

\subsection{Bandwidth selection}\label{bandwidth-selection}

For \(x\in \mathcal{S}\) and \(h>0\), the selection of the bandwidth parameter to be used with a given kernel \(K_{x,h}\) was performed by the cross-validation method in \CRANpkg{kernopt} package. Nevertheless, other procedures are available in the literature such as Bayesian bandwidth selection. We determined the bandwidth \(h_{CV}\) by minimizing in \(h\) the cross-validation function
\begin{equation*}
   CV(h)= \sum_{x\in\mathcal{S}}\left (\widehat{f}_{K,h}(x)\right )^{2} -\frac{2}{n}\sum_{i=1}^{n}\widehat{f}_{K,h,-i}(X_{i}),
\end{equation*}
where \(\widehat{f}_{K,h,-i}\) is calculated as \(\widehat{f}_{K,h}\) but by excluding the observation \(X_{i}\). The cross-validation procedure was implemented through the function \texttt{cv\_bandwidth} for the discrete kernels considered.

\section{Discrete optimal symmetric kernel}\label{sec3-optimal-kernel}

For \(x \in \mathcal{S}\) and \(h>0\), to determine a discrete symmetric ``optimal'' associated kernel \(K_{x,h}^{opt}(\cdot)\geq0\), we minimize the sum \(\sum_{y\in\mathcal{S}_{x}}K^{2}_{x,h}(y)\) from the expression of MISE in \eqref{eq:mise}. That consists in a minimization problem under the following constraints (see conditions in \eqref{eq:condition-symm-kern}):
\[\text{C}_{1}:\sum_{y\in\mathcal{S}_{x}}K_{x,h}(y)=1,  \text{C}_{2}:\sum_{y\in\mathcal{S}_{x}}yK_{x,h}(y)=x \text{\ and \ }\text{C}_{3}:\sum_{y\in\mathcal{S}_{x}}y^{2}K_{x,h}(y)=\tau^{2}, \]
which can be solved by using the method of Lagrange multiplier. A discrete expression of the Lagrange function is given by
\[\mathcal{L}(K)= \sum_{y\in\mathcal{S}_{x}}\left (K_{x,h}^{2}(y) + \lambda_{1} K_{x,h}(y)+ \lambda_{2}yK_{x,h}(y)+ \lambda_{3}y^{2}K_{x,h}(y)\right )+ \beta,\]
where the constant \(\beta\) does not depend on \(K_{x,h}\) and \(\lambda_{i}\) for \(i\in \{1,2,3\},\) are the Lagrange multipliers corresponding to constraints C\(_1\)-C\(_3\).

The discrete symmetric ``optimal'' kernel has the general form
\begin{equation*}\label{eq:Kh_opt}
K_{x,h}^{opt}(y,\lambda_{1},\lambda_{2},\lambda_{3})=-\frac{1}{2}(\lambda_{1} +\lambda_{2}y+\lambda_{3}y^{2}).  
\end{equation*}
Without loss of generality, we investigate discrete symmetric ``optimal'' kernels \(K_{k;x,h}^{opt}\) on the support \(\mathcal{S}_{x}=\{x-k,\ldots,x,\ldots,x\pm k\}\), with \(k\geq1\) a fixed integer. We obtain the following result when solving the minimisation problem described above.

\subsubsection{\texorpdfstring{\emph{Theorem}}{Theorem}}\label{theorem}

\emph{
For \(x\in \mathcal{S}\) and a fixed integer \(k\geq1\), the discrete symmetric ``optimal'' kernel on \(\mathcal{S}_{x}=\{x-k,\ldots,x,\ldots,x+k\}\) that minimises the mean integrate squared error of the estimator \(\widehat{f}_{K,h}\) in \eqref{eq:estim-fn} is given by
\begin{equation*}
 K_{k;x,h}^{opt}(y)
 =\lambda\bigg(\frac{3k^{2}+3k-1}{5} -(x-y)^{2} \bigg)+\frac{h}{2k+1}, y\in\mathcal{S}_{x},  
\end{equation*}
where \(\lambda=15(1-h)/\big((2k+1)(4k^{2}+4k-3)\big)>0\) and \(3/5(1-1/k)<h<1\).}\\

The bandwidth parameter \(h\) of the ``optimal'' kernels \(K_{k;x,h}^{opt}\) is bounded with a lower bound \(h_{lower}= 3/5(1-1/k)\geq0\), which particularly depends on \(k\geq 1\). The distribution of that kernel can be plotted as follows in the package (Figure \ref{fig:kernopt-optimal}):

\begin{verbatim}
x <- 5
z <- 0:10
h <- c(0.1, 0.4, 0.7, 0.9)
k <- 1

K_opt <- matrix(
  data = 0,
  nrow = length(z),
  ncol = length(h)
)

for (i in 1:length(h))
{
  K_opt[, i] <- discrete_kernel(kernel = "optimal", x, z, h[i], k)
}
plot(
  z,
  K_opt[, 1],
  xlab = "x",
  ylab = "Probability",
  ylim = c(0, 1),
  pch = 1
)
lines(z, K_opt[, 1], lty = 1)
for (i in 2:length(h))
{
  points(z, K_opt[, i], xlab = "z", pch = i)
  lines(z, K_opt[, i], lty = i)
}
legend(
  "topleft",
  c("h=0.1", "h=0.4", "h=0.7", "h=0.9"),
  lty = 1:4,
  pch = 1:4,
  cex = 1.6
)
\end{verbatim}

\begin{figure}
\centering
\pandocbounded{\includegraphics[keepaspectratio,alt={\label{fig:kernopt-optimal}Distribution of discrete symmetric `\,'optimal'\,' kernel (k=1) at the target point x=5 for various bandwidth parameters h.}]{RJ-2026-018_files/RJ-2026-018_files/figure-latex/kernopt-optimal-1.pdf}}
\caption{\label{fig:kernopt-optimal}Distribution of discrete symmetric `'optimal'\,' kernel (\(k=1\)) at the target point \(x=5\) for various bandwidth parameters \(h\).}
\end{figure}

We can formulate the following proposition on the limit distributions of \(K_{k;x,h}^{opt}\) as a function of \(h\).

\subsubsection{\texorpdfstring{\emph{Proposition 2}}{Proposition 2}}\label{proposition-2}

For \(x\in\mathcal{S}\), \(h\in(0,1)\) and integer \(k\geq1\), we consider the discrete symmetric ``optimal'' kernel \(K_{k;x,h}^{opt}(\cdot)\) with the support \(\mathcal{S}_{x}=\{x-k,\ldots,x,\ldots,x+k\}\).

\begin{enumerate}
\def\labelenumi{\roman{enumi}.}
\tightlist
\item
  \emph{When \(h\to 1\), the limit ``optimal'' kernel \(K_{k;x,1}^{opt}\) follows the discrete uniform distribution on the support \(\mathcal{S}_{x}\).}
\item
  \emph{When \(h\to 0\), the limit ``optimal'' kernel \(K_{k;x,0}^{opt}\) follows the Dirac distribution on the support \(\mathcal{S}_{x}\).}\\
\end{enumerate}

\paragraph{MISE of estimators as a function of discrete symmetric kernels.}

For the three discrete symmetric associated kernels considered, the MSSE of \(\widehat{f}_{K,h}\) in \eqref{eq:mise} can be also given by
\begin{eqnarray*}\label{eq_2:mise}
\mathrm{MSSE}\left (\widehat{f}_{K,h}(x)\right )
&=&\frac{1}{n}(\Pr(\mathcal{K}_{x,h}=x))^{2}C_{1}+ \frac{1}{4}(\text{Var}(\mathcal{K}_{x,h}) )^{2}C_{2} + o\bigg(\frac{1}{n} + h^{2} \bigg)\\
&=&\mathrm{AMSSE}\left (\widehat{f}_{K,h}(x)\right ) + o\bigg(\frac{1}{n} + h^{2} \bigg),\notag
\end{eqnarray*}\\
with \(C_{1}=\sum_{x\in\mathcal{S}}f(x)(1-f(x))>0\) and \(C_{2}=\sum_{x\in\mathcal{S}}(f^{(2)}(x))^{2}>0\).
For instance, the AMSSE of the estimator using the discrete symmetric ``optimal'' kernel is given by
\begin{eqnarray*}\label{eq:amise_Kopt}
\mathrm{AMSSE}\left (\widehat{f}_{K^{Opt},h}(x)\right )
&=&\frac{1}{n}\frac{\big(3(3k^2+3k-1) - 5hk(k+1)\big)^{2}}{(2k+1)^{2}(4k^{2}+4k-3)^{2}}C_{1}+ \frac{h^{2}k^{2}(k+1)^{2}}{36}C_{2}, 
\end{eqnarray*}
with \(\lambda=15(1-h)/((2k+1)(4k^{2}+4k-3))\).
Likewise, the AMSSE is given as follows, when considering the estimator using the discrete symmetric triangular kernel:
\begin{equation*}\label{eq:amise_Ktriang}
\mathrm{AMSSE}\left (\widehat{f}_{K^{Triang},h}(x)\right )=\frac{1}{n}\left(1-2hA(p)\right)^{2}C_{1}+ \frac{1}{4}h^{2}V^{2}(p)C_{2}
\end{equation*}
and the estimator using the discrete Epanechnikov kernel:
\[\mathrm{AMSSE}\left (\widehat{f}_{K^{Epan},h}(x)\right )=\frac{1}{n}\frac{9h^{2}}{(4h^{2}-1)^{2}}C_{1}+ \frac{1}{100}(h^{2}-1)^{2}C_{2}.\]
When considering discrete symmetric ``optimal'' and triangular kernels, an analytical expression of \(h>0\) that minimizes \(\mathrm{AMSSE}\) of \(\widehat{f}_{K,h}\) could be found, unlike Epanechnikov kernel for which only a numeric approximation was available.

\section{Package presentation on simulated data}\label{sec4-presentation-simu}

We illustrate the \CRANpkg{kernopt} package by applying it to simulated data. \citet{SengaDurrieu2024} performed a simulation study using different sample sizes \(n\) from various pmfs \(f\), i.e., a Poisson distribution, a geometric distribution and a negative binomial distribution. The code below presents simulations for estimating count data sets of size \(n=100\) that follow the Poisson distribution \({\cal P} (\lambda=2)\) (Figure \ref{fig:kernopt-simulated-data}). The kernel estimator was applied with each of three discrete symmetric kernels presented above in comparison with the count binomial asymmetric kernel that follows the binomial distribution \({\cal B}(x+1,(x+h)/(x+1))\) with support \({\cal S}_x=\{0,1,\ldots,x+1\}\), for \(x \in {\cal S} \subseteq {\mathbb N}\) and \(h \in (0,1]\) as seen in \citet{KokonSenga11}. The cross-validation method was used to select the bandwidth parameter \(h_{cv}\).

For a fixed kernel \(K\) and a bandwidth parameter \(h=h_{cv}>0\), the performance of the nonparametric kernel estimator \(\widehat{f}_{K,h_{cv}}\) was compared to the true pmf \(f\) of simulated data by using the Summed Squared Error (SSE) given by
\[\mathrm{SSE}(h_{cv})=\sum_{x\in\mathbb{N}}(\widehat{f}_{K,h_{cv}}(x) - f(x))^{2}. \]
In addition, {]}\{color=``blue''\} we can consider a goodness-of-fit test derived from the chi-squared statistic to test the null hypothesis \(H_0\): ``the true pmf \(f\) is distributed according to the pmf \(P_{0}\)'', for fixed \(\alpha=0.05\) for example. As the \(P_{0}\) is unknown, we replace \(P_{0}\) by a kernel estimator \(\widehat{f}_{K,h}\) for a fixed kernel \(K\) and a bandwidth parameter \(h=h_{cv}>0\). The statistic can be given by\\
\begin{equation*}
\chi _{0}^{2}=\sum_{x \in \{0,1,\ldots,N_{0}\}}\frac{\left(
n\widehat{f}_{K,h_{cv}}(x)- n_{i}\right)^{2}}{n\widehat{f}_{K,h_{cv}}(x)},
\end{equation*}
where \(n_{i}\) is the number of observations in each class \(x \in \{0,1,\ldots,N_{0}\}\). Since the distribution of the statistics test under \(H_0\) is unknown, we can use the bootstrap procedure to determine the distribution and the \(p\)-value. The Kolmogorov-Smirnov test can also be used for discrete data using the discrete kernel estimator \(\widehat{f}_{K,h_{cv}}\) of the pmf \(f\).

\begin{verbatim}
# Simulated data
#---------------
set.seed(18)
lambda <- 2
x <- 0:10
f <- dpois(x, lambda)
n <- 100
y <- sort(rpois(n, lambda))
data_sim <- as.data.frame(table(y))
f_sim <- c(f[1], f[2], f[3], f[4], f[5], sum(f[6:8])) # classes for chi-2 test

# Estimation with optimal kernel
#------------------------------
k <- 1
H <- seq((max(y) - min(y)) / 200, (max(y) - min(y)) / 2, length.out = 50)
hcv_opt_k1 <- cv_bandwidth(kernel = "optimal", y, H, k)
Fn_opt_k1 <- estim_kernel(kernel = "optimal", x, hcv_opt_k1, y, k)
SSE_opt_k1 <- sum((Fn_opt_k1 - f)^2)

# Estimation with triangular kernel
#---------------------------------
a <- 1
hcv_trg_a1 <- cv_bandwidth(kernel = "triang", y, H, a)
Fn_triang_a1 <- estim_kernel(kernel = "triang", x, hcv_trg_a1, y, a)
SSE_triang_a1 <- sum((Fn_triang_a1 - f)^2)

# Estimation with Epanechnikov kernel
#-----------------------------------
H <- seq(2, 10, 1)
hcv_epanech <- cv_bandwidth(kernel = "epanech", y, H, k = NULL)
Fn_epanech <- estim_kernel(kernel = "epanech", x, hcv_epanech, y, k = NULL)
SSE_epanech <- sum((Fn_epanech - f)^2)

# Estimation with Binomial kernel
#--------------------------------
H <- seq((max(y) - min(y)) / 500, 1, length.out = 50)
hcv_bin <- cv_bandwidth(kernel = "binomial", y, H)
Fn_bino <- estim_kernel(kernel = "binomial", x, hcv_bin, y, k = NULL)
SSE_bino <- sum((Fn_bino - f)^2)

# Frequency estimation
#---------------------
freq_sim <- as.numeric(as.character(data_sim[, 2])) # count data
f0_sim <- c(freq_sim / sum(freq_sim), rep(0, length(x) - length(freq_sim))) # frequency
SSE_0_sim <- sum((f0_sim - f)^2)

# Graph
#-----
plot(x, f,
  xlab = "x", ylab = "Frequency", xlim = c(0, 11), ylim = c(0, 0.35), type = "h",
  lwd = 2, col = "grey", main = ""
)
points(x + 0.15, f0_sim, type = "h", lty = 1, lwd = 2.75)
points(x + 0.3, Fn_opt_k1, type = "h", lty = 3, col = "grey", lwd = 2.75)
points(x + 0.45, Fn_epanech, type = "h", lty = 2, lwd = 2)
points(x + 0.6, Fn_triang_a1, type = "h", lty = 2, col = "grey", lwd = 2)
points(x + 0.75, Fn_bino, type = "h", lty = 3, lwd = 2)
legend("topright", c("Poisson f", "Frequency", "Optimal", "Epanechnikov", "Triangular", "Binomial"),
  lty = c(1, 1, 3, 2, 2, 3), lwd = 2, col = c("grey", "black", "grey", "black", "grey", "black"), inset = .0
)
\end{verbatim}

\begin{figure}
\centering
\pandocbounded{\includegraphics[keepaspectratio,alt={\label{fig:kernopt-simulated-data}Barplot of the frequency estimation of a count data set of sample size n=100 (black lines) simulated from a Poisson distribution f:=\{\textbackslash cal P\} (\textbackslash lambda=2) (grey lines), compared with estimations using discrete kernels (dotted lines). Summed Squared Error (SSE) was also calculated resulting in the following results as a function of the kernel used: Optimal k=1, SSE=0.0022 (h\_\{cv\}=0.95); Triangular a=1, SSE=0.0022 (h\_\{cv\}=3.22); Epanechnikov, SSE=0.0023 (h\_\{cv\}=2) ; Binomial, SSE=0.0069 (h\_\{cv\}=0.11); Frequency, SSE=0.0084.}]{RJ-2026-018_files/RJ-2026-018_files/figure-latex/kernopt-simulated-data-1.pdf}}
\caption{\label{fig:kernopt-simulated-data}Barplot of the frequency estimation of a count data set of sample size \(n=100\) (black lines) simulated from a Poisson distribution \(f:={\cal P} (\lambda=2)\) (grey lines), compared with estimations using discrete kernels (dotted lines). Summed Squared Error (SSE) was also calculated resulting in the following results as a function of the kernel used: Optimal k=1, SSE=\(0.0022\) (\(h_{cv}=0.95\)); Triangular a=1, SSE=\(0.0022\) (\(h_{cv}=3.22\)); Epanechnikov, SSE=\(0.0023\) (\(h_{cv}=2\)) ; Binomial, SSE=\(0.0069\) (\(h_{cv}=0.11\)); Frequency, SSE=\(0.0084\).}
\end{figure}

The choice of the bandwidth parameter \(h\) is crucial, as usual in discrete and continuous nonparametric kernel estimations. Nevertheless, the quality of the fit is more sensitive to the kernel chosen for discrete kernel estimates \(K\) than for continuous kernel estimates.

Broader comparisons of the discrete symmetric ``optimal'' kernel with discrete symmetric and asymmetric kernels, and the frequency estimator, were carried out on \(N=100\) replicates of different samples sizes \(n\in\{25,100,250,500\}\) (i.e., small, medium and large) in \citet{SengaDurrieu2024}, to strengthen evidence for the advantages of the ``optimal'' kernel, in particular for the small or moderate sample sizes \(n\in\{25,100\}\).

\section{Application on real data}\label{sec5-application-real-data}

SIMTAP project (\url{https://www.simtap.eu}) aimed to develop sustainable aquaculture production production system that, in particular, contributes to reduce fish feed inputs and resources consumption. We used data from an experiment in which gilthead seabream (sparus aurata) were stocked in 1.6 m\(^{3}\) tanks at a density of 1.5 kg\(\cdot\)m\(^{-3}\). Fish were reared for 46 days in a single recirculating aquaculture system composed of three rearing tanks. At the beginning of the experiment, a number \(n=200\) of the fish were individually weighed (dg), and their length at the caudal fork (mm) was measured. The nonparametric kernel estimator was applied to provide smoothed count distributions of weight and length of fish measured. As the true pmf \(f\) of real data was unknown, the performance of the nonparametric kernel estimator \(\widehat{f}_{K, h}\) was compared to that of the frequency estimator \(f_{0}\) by using {SSE previously defined in simulations. In addition, the cross-validation method was used to select the bandwidth parameter}.

The following code illustrates this application on real experimental fish data with the \CRANpkg{kernopt} package. The corresponding data are included in the package as an example data set and can be load through the \emph{utils} library \texttt{data()} function as ``fish\_data'' data set, e.g., \texttt{data(\textquotesingle{}fish\_data\textquotesingle{},\ package\ =\ \textquotesingle{}kernopt\textquotesingle{})}.
It also uses the \CRANpkg{kernopt} \texttt{cv\_bandwidth()} function to compute the estimate an optimal bandwidth parameter from the cross-validation procedure and then the \texttt{estim\_kernel()} function in order to calculate the corresponding SSE. Empirical frequencies and discrete kernel estimates of count data of weight of fish were graphically compared (Figure 5).

\begin{verbatim}
#### Data
data("fish_data", package = "kernopt")
y <- fish_data$weight
data1 <- as.data.frame(table(y))
x <- as.numeric(as.character(data1[, 1]))
freq_weight <- as.numeric(as.character(data1[, 2]))
f0 <- freq_weight / sum(freq_weight)
n <- sum(freq_weight)

#### Optimal kernel
H <- seq((max(y) - min(y)) / 200, (max(y) - min(y)) / 2, length.out = 50)
k <- 1
hcv_opt_k1 <- cv_bandwidth(kernel = "optimal", y, H, k = 1)
fn_opt_k1 <- estim_kernel(kernel = "optimal", x, hcv_opt_k1, y, k = 1)
SSE_opt_k1 <- sum((fn_opt_k1 - f0)^2)

#### Epanechnikov
H <- seq(2, 10, 1)
hcv_epanech <- cv_bandwidth(kernel = "epanech", y, H)
fn_epanech <- estim_kernel(kernel = "epanech", x, hcv_epanech, y, k = NULL)
SSE_epanech <- sum((fn_epanech - f0)^2)

# Graph
par(mfrow = c(1, 2))
plot(x, f0,
  col = "black", axes = F, lwd = 3, ylab = "", xlab = "Weight (dg)",
  ylim = c(0, 0.06), xlim = c(41, 132), type = "h",
  main = "(a)", cex.axis = 1, cex.lab = 1
)
axis(1, at = x, cex.axis = 1, cex.lab = 1)
axis(2)
box()
points(x + 0.4, fn_opt_k1, lwd = 3, col = "grey", lty = 1, type = "h")
plot(x, f0,
  col = "black", axes = F, lwd = 3, xlab = "Weight (dg)", ylab = "",
  ylim = c(0, 0.06), xlim = c(41, 132), type = "h",
  main = "(b)", cex.axis = 1, cex.lab = 1
)
axis(1, at = x, cex.axis = 1, cex.lab = 1)
axis(2)
box()
points(x + 0.4, fn_epanech, lwd = 3, col = "grey", lty = 1, type = "h")
\end{verbatim}

\begin{figure}
\centering
\pandocbounded{\includegraphics[keepaspectratio,alt={\label{fig:kernopt-real-data}Density estimations (gray lines) of count distributions of weight (dg) (black lines) of fish by using optimal and Epanechnikov discrete kernels with bandwidth parameter (h\_\{cv\}) from the cross-validation procedure. Summed Squared Error (SSE) was also calculated. (a) Optimal k=1: SSE=0.0006 (h\_\{cv\}=0.47); (b) Epanechnikov: SSE=0.0033 (h\_\{cv\}=9).}]{RJ-2026-018_files/RJ-2026-018_files/figure-latex/kernopt-real-data-1.pdf}}
\caption{\label{fig:kernopt-real-data}Density estimations (gray lines) of count distributions of weight (dg) (black lines) of fish by using optimal and Epanechnikov discrete kernels with bandwidth parameter (\(h_{cv}\)) from the cross-validation procedure. Summed Squared Error (SSE) was also calculated. (a) Optimal k=1: SSE=\(0.0006\) (\(h_{cv}=0.47\)); (b) Epanechnikov: SSE=\(0.0033\) (\(h_{cv}=9\)).}
\end{figure}

\subsection{Acknowledgments}\label{acknowledgments}

We thank the \href{https://www.simtap.eu}{SIMTAP} project (Self-sufficient Integrated Multi-Trophic AquaPonic systems) supported by the PRIMA program (Partnership for Research and Innovation in the Mediterranean Area) under grant (18110-2) for providing the data.

\bibliography{RJreferences.bib}

\address{%
Gilles Durrieu\\
Univ Bretagne Sud, LMBA CNRS UMR 6205\\%
F-56000 Vannes, France\\
%
%
\textit{ORCiD: \href{https://orcid.org/0000-0003-0375-912X}{0000-0003-0375-912X}}\\%
\href{mailto:gilles.durrieu@univ-ubs.fr}{\nolinkurl{gilles.durrieu@univ-ubs.fr}}%
}

\address{%
Thomas Fillon\\
Univ Bretagne Sud, LMBA CNRS UMR 6205 \& IRISA CNRS UMR 6074\\%
F-56000 Vannes, France\\
%
%
\textit{ORCiD: \href{https://orcid.org/0000-0003-1207-8414}{0000-0003-1207-8414}}\\%
\href{mailto:thoma.fillon@univ-ubs.fr}{\nolinkurl{thoma.fillon@univ-ubs.fr}}%
}

\address{%
Tristan Senga Kiessé\\
INRAE, Institut Agro \& UMR 1069 SAS\\%
F-35000 Rennes, France\\
%
%
\textit{ORCiD: \href{https://orcid.org/0000-0003-2710-5825}{0000-0003-2710-5825}}\\%
\href{mailto:tristan.senga-kiesse@inrae.fr}{\nolinkurl{tristan.senga-kiesse@inrae.fr}}%
}
